\itemhead{p.103 练习 · 第2题（教材p.103）}{XJ-B1-C04-PR-P103-Y1012}
\label{XJ-B1-C04-PR-P103-Y1012}
{\footnotesize 湘教版教材p.103 / PDF第110页 / p.103 练习 / 第2题}
\field{原题}{根据幂函数 $y=x^\alpha$ 的图象填空：
\subq{1}当 $\alpha>1$，$x\ge0$ 时，图象恒过点\blank 和\blank；当 $0<x<1$ 时，图象在直线 $y=x$ 的\blank 方；当 $x>1$ 时，图象在直线 $y=x$ 的\blank 方。
\subq{2}当 $0<\alpha<1$，$x\ge0$ 时，图象恒过点\blank 和\blank；当 $0<x<1$ 时，图象在直线 $y=x$ 的\blank 方；当 $x>1$ 时，图象在直线 $y=x$ 的\blank 方。
\subq{3}当 $\alpha<0$，$x>0$ 时，图象恒过点\blank。}
\field{参考答案}{（1）$(0,0),(1,1)$；下、上。（2）$(0,0),(1,1)$；上、下。（3）$(1,1)$。}
{\small \field{关联知识}{P03 幂函数的概念、图象与性质（教材帮p.125、126、127 / PDF134、135、136）；F04 函数三种表示法、图象与读数（教材帮p.85、86、93、94 / PDF94、95、102、103）}
\field{方法与易错点}{设 $f(x)=x^\alpha$，用已知点求指数，再查该指数的定义域、奇偶性和各单调区间。 易错：把 $3x^2$ 当标准幂函数；将第一象限性质直接用于负半轴；分数负指数遗漏0禁点。}
\field{答案依据}{答案册相关题组：PDF第30、31页（册内页29、30）；按当前题库原题逐项重算，不按旧版题号机械对套。}}
